\section{Discussion}

This study represents a comprehensive, population-wide evaluation of one-year stroke outcomes in England using a dataset of 388,345 stroke patients linked across national clinical registries, hospital administrative databases, and vaccination records within the NHS England Secure Data Environment (SDE). By utilizing the patient-centred metric of \textit{home-time} (days alive and outside of NHS inpatient care) and matching stroke survivors with age-, sex-, and ethnicity-matched controls ($n = 236,630$), we have quantified the real-world longitudinal burden of stroke with added statistical resolution. Furthermore, by supplementing a classic Cox proportional hazards model with explainable machine learning (XGBoost and SHAP), this study uncovers complex non-linear relationships, identifies the protective effects of COVID-19 vaccination, and highlights ethnic differences in long-term stroke survival.

Our use of home-time (days alive and outside NHS inpatient care) is a  pragmatic, registry-derived, patient-centred outcome that may be collected at large scale. Our findings complement validation work by Sung et al. \cite{sung2020home}, who demonstrated that post-stroke home-time strongly correlates inversely with 1-year mRS scores (Spearman's $\rho$ ranging from $-0.69$ to $-0.88$), proving that home-time functions as a highly reliable surrogate for functional outcome. While Sung et al. validated this tool within a localized registry framework, our study scales this metric to a large scale national electronic health record ecosystem. We demonstrate that stroke patients lose an average of 86 days at home, die 66 days sooner, and spend 21 more days in inpatient NHS care in their first year compared to demographically identical peers, providing a clear concrete metric of the absolute burden of stroke. The results reported here  supports the case for home-time as a potential standard secondary endpoint in future national stroke audits.

The one-year fatality rate of 22.0\% after stroke is similar to that found in other studies. It is consistent with the downward trajectory in general-practice cohort fatality (from 29.2\% to 22.2\% in men between 1997 and 2005) \cite{gulliford_declining_2010} and with the South London Stroke Register's reduction in one-year adjusted case-fatality to 20.4\% by 2012--2015 \cite{wafa_long-term_2020}, and is similar to the 18.2\% one-year fatality reported in a large Swedish cohort \cite{ullberg_changes_2015}.

One-year survival was negatively affected by increased age, higher pre-stroke disability, and greater stroke severity, results consistent with previous findings \cite{nambiar2022}. One of our most striking findings is that Black and Asian stroke patients experienced a substantial reduction in 1-year mortality risk (65\% and 67\% hazard ratios compared to White patients, respectively), while SHAP models simultaneously predicted higher baseline time in inpatient care for Black, Mixed, and Other ethnic minorities, while also predicting higher survival for non-white patients. This distinct survival advantage closely mirrors historical data from the South London Stroke Register (SLSR); Wang et al. \cite{wang2013trends} analyzed a 16-year multi-ethnic cohort and found that Black Caribbean and Black African patients had a significantly reduced risk of all-cause mortality compared to White patients (adjusted HR 0.85 and 0.61, respectively), particularly among older individuals. Wang et al. hypothesized that this paradox might be driven by a `healthy migrant effect' or differences in acute stroke care pathways. Our study shows that this survival advantage is not merely a regional anomaly localized to South London but is a persistent nationwide phenomenon across England that remains significant even after rigorous adjustment for baseline stroke severity, age, clinical subtypes, and socioeconomic deprivation. However, it remains a possibility that these results are an artifact of data recording - results based on registries are for those that have attended hospital and being diagnosed as stroke; if there is a delay in decision to attend hospital - death may be occurring more before arrival at hospital, discounting the measured risk of death after arrival at hospital. The largest and most recent large study, using the American Heart Association's \textit{Get With The Guidelines-Stroke} registry of over 600,000 patients presenting between mid-2015 and 2019, found that Black race, in the US, was independently associated with a prolonged symptom-onset-to-ED-arrival interval of an additional 28 minutes compared with White patients \cite{royan2024disparities}. The concurrent increase in inpatient care days observed for non-White ethnicty in our study suggests that ethnic variations warrant deeper exploration. 

Explainable machine learning offers a flexible approach to predicting outcomes after stroke \cite{pearn2026machine}. Machine learning allow for non-linear effects, and for all possible interactions between patient characteristics, and the addition of explainability with SHAP allows for analysis and visualisation of these effects. The use of a cut-off of 12 months survival, and access to independent death data should avoid the '\textit{competing risks}' problem that can exist with models where patients might be lost to follow-up due to death. It is important to acknowledge that our general models are predictive rather than causative. We can still observe clear predictive patterns. Stroke severity, not surprisingly, has a non-linear relationship with time in NHS care - patients with moderate and moderate-severe stroke severity have the highest time in NHS care in the year after stroke, after adjusting for other factors. The lowest stroke servility patients have lower time in care, and most time at home, whereas the highest stroke severity patients have lower time in care, and most time time dead in the year after stroke. Prior disability has a similar pattern - medium pre-stroke disability have the the longest time in NHS care after adjusting for other factors, whereas low pre-stroke disability predicts shorter time in care and longer time at home, whereas high pre-stroke disability predicts shorter time in care and longer time dead in the year after stroke. Age at stroke appears to have little association time in NHS care, but less time at home and more time dead after stroke. Ethnicity, as commented upon shows a pattern where non-White patients have longer time in NHS care, longer time dead, and less time at home in the year after stroke. After adjusting for other factors, deprivation  appears to have little association with time in care, at home, or dead. A recent Covid infection, without vaccination, was associated with slightly longer time in NHS care, less time at home, and more time dead. This was very largely ameliorated by Covid vaccination. Outcomes after Covid vaccination alone (without recorded Covid infection) were very similar to people with no recorded Covid infection. The results of COVID and vaccination are consistent with reported results \cite{el2025vaccination, rizzo2022covid}.

Finally, our analysis of the reperfusion cohort ($n = 136,395$), which best isolates the effect of treatment, strongly reinforces the "\textit{time is brain}" paradigm \cite{saver2006time} on a national scale. Early dual therapy (IVT and MT under 270 minutes) yielded the maximum optimization of home-time and reduction in mortality, while reducing time in NHS care. Conversely, late mechanical thrombectomy alone (beyond 4.5 hours) was associated with an increased NHS inpatient time, reduced home-time, and higher mortality risks. In this study we used a single division of time of 4.5 hours from stroke onset (the window of licensed use for thrombolysis). These results suggest that a deeper investigation, with more detailed analyses of time-to-treatment, is warranted.

\subsection{Strengths and limitations}

The primary strength of this study is the combination of multiple national data sets, to allow a very large scale study of predictors of inpatient use, home-time, and death after stroke. We were able to compare this cohort with a control population matched on age,sex, and ethnicity.

Several limitations should be acknowledged:

\begin{itemize}
    \item \textit{Definition of `home-time'}:  Because home-time was calculated as days spent alive outside of NHS inpatient beds, patients discharged to private, self-funded nursing homes or long-term social care facilities may be classified as being \textit{at home}. Consequently, our metric represents an index of independence from acute/sub-acute NHS hospital structures rather than absolute residential autonomy.
    \item \textit{Potential residual confounding}: Despite adjusting for a wide array of clinical features, it is possible that unmeasured factors have introduced bias into the results.
\end{itemize}

\subsection{Conclusion}

By leveraging population-wide data, traditional statistical methods, and  explainable machine learning, this study establishes that stroke in England inflicts a profound, quantifiable deficit on a patient's first year of post-event life, resulting in an average loss of 86.4 home days compared to matched controls. Survival and home-time are highly dependent on hyper-acute treatment timing, prior disability, and age. Crucially, our models reveal that COVID-19 vaccination effectively buffers vulnerable stroke populations against the severe long-term risks of infection, and we confirm a striking nationwide survival advantage among Black and Asian stroke survivors. Our results have emphasised the continual need to focus on speed of treatment, in a real world population. 